\documentclass{article}

% Preamble
\usepackage[margin=1in]{geometry}
\usepackage{amsmath}
\usepackage{hyperref}
\usepackage{enumitem}

\title{AgentWhetters: Spatial Reasoning Skills for Cost-Effective\\Block-Building Agents}
\author{
    Whitten, Paul\\
  \texttt{pcw@case.edu}
  \and
  Chen, Li-Jen\\
  \texttt{chenlijen@gmail.com}
  \and
  Baddam, Sharath\\
  \texttt{sbaddam@alumni.purdue.edu}
}
\date{March 2026}

\begin{document}

\maketitle

\begin{abstract}
We present \textbf{AgentWhetters BWIM}, a purple agent for the AgentBeats
\emph{Build What I Mean} (BWIM) benchmark that constructs block structures on a
$9\!\times\!5\!\times\!9$ grid from sometimes underspecified natural-language
instructions. The system originated as an investigation into enabling small,
cost-efficient language models, originally targeting edge deployment with
limited compute, to perform reliable 3D spatial reasoning.  It has since been
adapted as a competition entry for the AgentX AgentBeats Phase~2 Sprint~1, where
it is evaluated on structural accuracy and question efficiency against both a
rational and an unreliable architect partner.  Rather than relying on large
frontier models, our approach augments GPT-4o-mini with a modular skills
pipeline that offloads spatial computation to deterministic code, reducing the
burden on the language model to tasks it handles well: instruction
comprehension, ambiguity detection, and structured plan generation.  The name
\emph{AgentWhetters} reflects our view that competitions such as AgentX
AgentBeats offer a valuable opportunity to \emph{whet} or \emph{sharpen} our
skills in developing, evaluating, and iterating on autonomous agents under
realistic constraints.
\end{abstract}

\section{Architecture}

The agent decomposes each building instruction through a five-stage pipeline:
\textbf{parse}, \textbf{plan}, \textbf{detect underspecification},
\textbf{execute}, and \textbf{format}.  If any stage fails, control falls back
to a direct LLM call with a carefully engineered system prompt.

\begin{enumerate}[leftmargin=*]
  \item \textbf{Instruction Parser.} Extracts the green agent's instruction
    text, separates the building directive from any existing grid state
    (START\_STRUCTURE), and normalizes the representation for downstream stages.

  \item \textbf{Structure Analyzer.} Before planning, the current grid state is
    analyzed to detect geometric primitives rows, stacks, L-shapes,
    T-shapes and produce a structured description.  This pre-computed geometry
    is injected into the planner prompt so the language model reasons over named
    shapes rather than raw coordinate lists.

  \item \textbf{Build Planner (LLM).} The language model decomposes the
    instruction into a sequence of atomic build steps expressed as typed JSON
    actions (\texttt{stack}, \texttt{place\_relative}, \texttt{extend\_row},
    etc.), each with a color, count, and position specification.  Nine worked
    examples in the system prompt covering: chains, L-shapes, T-shapes, and
    edge placements anchor the model's spatial reasoning within the coordinate
    system.

  \item \textbf{Plan Verifier and Deterministic Corrections.} A rule-based
    post-planner validates the LLM-generated plan against the instruction text
    and the current grid geometry. Four deterministic correction passes run
    before execution:
    \begin{enumerate}
      \item \emph{Direction consistency:} compares directional words in the
        instruction (``left,'' ``right,'' ``front,'' ``behind'') against the
        plan's \texttt{direction} fields and flips any contradiction.
      \item \emph{Each-end cap correction:} when the instruction contains
        ``each end'' or ``both ends'' and the plan extends a row, the system
        recomputes the new endpoints from the extension geometry and relocates
        any cap-placement steps to the correct positions.
      \item \emph{T-shape extend correction:} runs the structure analyzer on
        the starting grid to detect T-shaped geometry, computes the stem's axis
        and growth direction from block coordinates, and replaces invalid or
        incorrect extend directions (the LLM frequently generates
        \texttt{on\_top}, which is not a valid horizontal direction).
      \item \emph{Stacking and count plausibility:} checks
        horizontal-versus-vertical intent and block-count consistency.
    \end{enumerate}
    These corrections operate on general geometric properties of the grid and
    instruction text. They do not reference specific stimuli or expected
    outputs and generalize to arbitrary instructions and grid configurations.

  \item \textbf{Spatial Executor.} A deterministic engine executes each plan
    step on an in-memory grid model, resolving relative references
    (``leftmost,'' ``the red one''), computing gravity-based y-coordinates, and
    chaining positional context across steps.  No LLM call is involved at this
    stage; all coordinate arithmetic is exact.  A same-color skip-forward rule
    in the \texttt{extend\_row} handler detects when the specified start
    position is already occupied by a block of the same color and advances the
    start by one grid step in the extension direction before placing.  This
    prevents vertical stacking when the LLM intends horizontal extension from
    an existing block.

    A key design choice is that the LLM planner operates in what is effectively
    a two-dimensional problem space.  Each plan step specifies only a horizontal
    position $(x, z)$ and an action type; the executor computes the vertical
    coordinate automatically by scanning the column for existing blocks and
    placing each new block at the next available height, much like dropping a
    piece in Connect~4.  This means the language model never needs to perform
    $y$-coordinate arithmetic (ground level at $y\!=\!50$, each block adding
    $+100$) or track how many blocks already occupy a column.  In our initial
    approach, the LLM was asked to produce the complete
    \texttt{[BUILD];Color,x,y,z;\ldots} response directly, which required it
    to enumerate every block's three-dimensional coordinates and to correctly
    chain height calculations across multiple stacking operations.  That
    approach suffered from frequent off-by-one $y$-errors, duplicate blocks at
    the same position, and miscounted stack heights.  The current executor
    eliminates these failure modes entirely by confining the LLM to the
    two-dimensional layout it can reason about reliably and delegating all
    vertical placement to deterministic code.

  \item \textbf{Response Formatter.} Translates the final grid state into the
    protocol-required \texttt{[BUILD]} format and validates the output before
    transmission.
\end{enumerate}

\section{Underspecification Detection}

A critical challenge in the benchmark is that many instructions omit color or
block count.  Our agent employs a two-tiered detection strategy:

\begin{itemize}[leftmargin=*]
  \item \textbf{Missing color.} Heuristic analysis determines whether a color
    can be inferred from context (e.g., reusing the sole color mentioned) or is
    genuinely ambiguous.  In the latter case, the agent issues an \texttt{[ASK]}
    query to the architect.  Because a correct answer yields $+10$ while asking
    costs $-5$, the expected value of asking ($+5$) exceeds guessing ($\approx
    0$).  A multi-color disambiguation loop handles instructions that leave
    multiple colors unspecified.

  \item \textbf{Missing count.} Empirical analysis of the stimulus data
    revealed that heuristic count inference (copying an adjacent stack's height
    or defaulting to three) achieves only ${\sim}65\%$ accuracy, well below
    the $75\%$ threshold at which guessing dominates asking under the scoring
    function.  At $p = 0.65$, $\text{EV}_{\text{guess}} = 20(0.65) - 10 =
    +3.0$, while $\text{EV}_{\text{ask}} = +5.0$ per round, a gain of $+2.0$
    per count-underspecified round.

    A complication is that the architect answering questions is itself an LLM
    (GPT-4o-mini) that receives only the target structure as coordinates, not
    the original instruction.  When presented with a generic question such as
    ``How many blocks should be in the unspecified stack?'' the architect
    cannot determine which stack is being referenced and defaults to answering
    ``3 blocks'' regardless of the true count, producing a 23\% error rate in
    our experiments.  To address this, we generate \emph{color-specific}
    questions that name the color of the uncounted phrase, for example ``How
    many blue blocks should be in the blue stack?''  This allows the architect
    to identify the correct stack by counting blocks of that color in its
    coordinate data.  The answered count is then patched directly into the
    instruction text before the LLM planner runs, so the planner receives a
    fully specified instruction such as ``build a stack of 3 blue blocks''
    rather than the ambiguous ``build a blue stack.''

    Because only one question is permitted per round, the agent prioritizes
    color over count.  When multiple phrases lack counts, the first uncounted
    phrase with a known color is selected for the question and the remainder
    fall back to heuristic inference.  The heuristic resolves unspecified
    counts through a three-level cascade: first, it copies the height of an
    adjacent stack at the target position; second, if no neighbor exists, it
    uses the tallest stack currently on the grid; third, it defaults to three,
    which matches the modal count in the stimulus data (53\% of
    count-underspecified trials).
\end{itemize}

\section{Adaptive Prompt Enrichment}

We observe that certain spatial concepts such as``each end,'' chain references,
L-shape extensions, T-shape junctions are disproportionately error-prone for
small models.  Our \emph{adaptive prompt enrichment} system scans each incoming
instruction against a library of 15 pattern-matching rules.  When a rule fires,
a short concept definition with a few-shot worked example (concrete coordinates
showing correct and incorrect placements) is injected into the user prompt
immediately before the LLM call.  This places the relevant spatial rule and a
grounding example in the model's local attention window, adjacent to the
instruction it must interpret.  Rules are independent and composable; multiple
may fire on a single instruction without conflict.  The enrichment library is
validated against the full stimulus set to ensure zero false positives on
fully-specified instructions.

To guard against teaching the test, all few-shot examples are constructed with
different colors, counts, positions, and phrasing from the evaluation scenarios.
We measure unigram+bigram cosine similarity between each example and every
stimulus instruction.  The highest similarity across all examples is 0.71,
while real stimulus instructions reach pairwise similarities of up to 0.80
among themselves.  Our examples therefore share less surface overlap with any
individual test item than the test items share with each other.

\section{Additional Design Considerations}

\begin{itemize}[leftmargin=*]
  \item \textbf{Chain reference patching.} A dedicated pre-processing pass
    resolves chain references (``to the right of the green one'') that span
    multiple build steps, ensuring positional context propagates correctly.

  \item \textbf{Ask-once guard.} The agent tracks whether it has already asked a
    question in the current round, preventing redundant queries that would incur
    additional point penalties.

  \item \textbf{Cost efficiency.} The entire pipeline uses GPT-4o-mini, one of
    the lowest-cost commercial models available.  Deterministic execution of
    spatial computation means the LLM is called at most twice per round (planning
    plus at most one clarification re-plan), keeping token usage minimal.

  \item \textbf{Graceful degradation.} If the skills pipeline encounters an
    unrecoverable error at any stage, the agent falls back to a direct LLM call
    with a comprehensive system prompt, ensuring every round produces a valid
    \texttt{[BUILD]} response.
\end{itemize}

\section{Results}

Table~\ref{tab:results} summarizes representative runs across development
iterations. In local evaluation against the benchmark's 160-round scenario
(combining rational and unreliable architect partners), the agent's best run
achieves 95.6\% structural accuracy (153 of 160 rounds correct) with a
cumulative score of $+980$.

\begin{table}[h]
\centering
\small
\begin{tabular}{llcccc}
\hline
\textbf{Run} & \textbf{Changes} & \textbf{Score} & \textbf{Accuracy}
  & \textbf{T-shape} & \textbf{Each-end} \\
\hline
Baseline   & No count-ask        & $+780$ & 81.9\% & --  & --  \\
+Count-ask & Count questions      & $+720$ & 87.5\% & 0/4 & 0/4 \\
+Enrichment & Prompt enrichment   & $+780$ & 88.1\% & 0/4 & 1/4 \\
+Each-end fix & Deterministic cap fix & $+820$ & 89.4\% & 1/4 & 4/4 \\
+Example rewrite & Decontaminated examples & $+860$ & 91.9\% & 2/4 & 2/4 \\
+Extend skip & Executor same-color fix & $+980$ & 95.6\% & 4/4 & 4/4 \\
\hline
\end{tabular}
\caption{Score progression across development iterations (160 rounds each).}
\label{tab:results}
\end{table}

Adding count-specific \texttt{[ASK]} questions raised structural accuracy from
82\% to 87\% at the cost of a lower raw score due to question penalties.
Adaptive prompt enrichment and deterministic post-plan corrections recovered
the score while further improving accuracy. The each-end deterministic fix
alone moved that sub-task from 1/4 to 4/4 correct, contributing $+40$ net
points. A subsequent fix to the underspecification detector's question
targeting (asking about the correct uncounted stack rather than a
misidentified one) addresses 6 additional failure rounds in testing.
Rewriting few-shot examples to avoid stimulus contamination and adding a
fallback each-end correction path brought the overall accuracy to 91.9\%.

The final improvement, a same-color skip-forward rule in the spatial
executor, corrected a class of failures where \texttt{extend\_row} stacked
vertically on an existing block of the same color instead of extending
horizontally past it. This fix raised the score from $+860$ to $+980$ and
moved both T-shape (2/4 to 4/4) and each-end (2/4 to 4/4) sub-tasks to
perfect accuracy.

The 7 remaining failures in the best run fall into two categories.
First, 5 rounds where the green agent (architect) provided incorrect color
or count information in its instruction, for example naming a block color as
Green when the target structure contained Blue, or reporting 4 blocks as the
count for a stack that should have 3. These errors originate in the architect
LLM and cannot be corrected by the builder agent. Second, 2 rounds of LLM
spatial reasoning errors: one incorrect chain coordinate reference and one
L-shape with combined wrong-color and vertical stacking.

The dominant remaining error source is therefore the green agent, not the
purple agent's pipeline. Further accuracy gains would require either
architect-side improvements or cross-validation heuristics to detect
and compensate for architect errors.

\end{document}